\documentclass[11pt, letterpaper]{report}

%PAQUETES A UTILIZAR
\usepackage{helvet} %PAQUETE PARA UTILIZAR LA FUENTE HELVETICA
\usepackage{graphicx} %PAQUETE QUE PERMITE LA MANIPULACIÓN DE IMAGENES
\usepackage[top=2.5cm, bottom=2.5cm, left=3cm, right=3cm]{geometry} %PAQUETE PARA MODIFICAR LOS MÁRGENES DE LA HOJA
\usepackage{fancyhdr} %PAQUETE PARA ENCABEZADOS O PIES DE PÁGINA
\usepackage{apacite} %PAQUETE QUE OFRECE EL SOPORTE PARA CITAS APA
\usepackage{parskip} %PAQUETE PARA ESPACIADO AUTOMÁTICO ENTRE PÁRRAFOS
\usepackage[compact]{titlesec} %Paquete para eliminar los espacios en blanco entre títulos, subtítulos, etc
\usepackage{setspace} %PAQUETE PARA CAMBIAR EL INTERLINEADO
\usepackage{afterpage}
%\usepackage[Lenny]{fncychap} %PAQUETE PARA CAMBIAR EL ESTILO DE TÍTULOS DE CAPÍTULOS, SECCIONES Y SUBSECCIONES
\usepackage{enumitem} %PAQUETE PARA LA CONFIGURACIÓN DE LISTAS
\usepackage{listings} %PAQUETE PARA MODIFICAR EL COMPORTAMIENTO DE LOS LISTADOS
\usepackage[bookmarks=true, pdfstartview=FitH]{hyperref} %PAQUETE PARA AÑADIR VÍNCULOS A LA TABLA DE CONTENIDOS
\usepackage{config/docstyle} %PAQUETE CUSTOM PARA LA INCRUSTACIÓN DE DATOS PERSONALES
%\usepackage{showframe} %PAQUETE QUE PERMITE VISUALIZAR EL MARGEN DEL DOCUMENTO

%AQUÍ COMIENZA EL DOCUMENTO
\begin{document}
    %Referencias renombrado
\renewcommand{\bibname}{Bibliografía}

%Configurando parskip
\setlength{\parindent}{0cm}
\setlength{\parskip}{5mm}

%Configurando el espaciado entre "itemize/enumerate items"
\setlist[itemize]{noitemsep, topsep=3mm}
\setlist[enumerate]{noitemsep, topsep=3mm}

%Espacios entre autores en "Bibliografía"
\let\OLDthebibliography\thebibliography
\renewcommand\thebibliography[1]{
    \OLDthebibliography{#1}
    \setlength{\parskip}{3mm}
    \setlength{\itemsep}{3mm plus 0.3ex}
}
    \begin{titlepage}
    \begin{center}
        \begin{minipage}{0.25\textwidth} %0.15 para que quede todo en una linea
            \begin{flushleft}
                \includegraphics[scale=0.08]{img/logo-sep-vertical.png} %scale=0.0610
            \end{flushleft}
        \end{minipage}
        \hfill
        \begin{minipage}{0.4\textwidth} %0.6 para que quede todo en una linea
            \begin{center}
                \MakeUppercase{\nombreUniversidad}
            \end{center}
        \end{minipage}
        \hfill
        \begin{minipage}{0.25\textwidth} %0.15 para que quede todo en una linea
            \begin{flushright}
                \includegraphics[scale=0.045]{img/logo-itc.png} %scale=0.04
            \end{flushright}
        \end{minipage}
        
        \vspace{0.8cm}
        \MakeUppercase{Informe técnico de residencias profesionales}

        \nombreProyecto{}

        \vspace{0.8cm}
        \fechaInicioResidencia{} \textendash{} \fechaFinResidencia{}

        \nombreEmpresa{}

        \vspace{0.8cm}
        Presenta:

        \nombreAutor{}

        \numeroControl{}

        \especialidad{}

        \vspace{0.8cm}
        Asesor Externo

        \nombreAsesorExterno{}

        \vspace{0.8cm}
        Asesor Interno

        \nombreAsesorInterno{}

        \vspace{0.8cm}
        \fechaPresentacion{}
    \end{center}
\end{titlepage}

    \pagenumbering{Roman}
    \tableofcontents

    \newpage
    \pagenumbering{arabic}
    \chapter{Justificación}
    Mezfer es una empresa mexicana dedicada a solucionar los problemas del campo mediante la elaboración de productos dirigidos al sector agroalimentario. El creciente número de clientes llevó a Mezfer a crear ``Mezfer Rewards'', un programa de recompensas mediante el cual los clientes se inscriben a los diferentes `Programas de lealtad' para canjear puntos que reciben con la compra de productos.
    
    Actualmente, Mezfer Rewards es gestionado mediante 2 aplicaciones diferentes: 
    \begin{itemize}
        \item Una aplicación web para que los administradores puedan realizar la gestión del programa.
        \item Una aplicación móvil para que los clientes puedan realizar todo el proceso para obtener recompensas.
    \end{itemize}

    Ahora, Mezfer busca que tanto clientes como administradores puedan realizar sus respectivas actividades en una misma plataforma para así mantener un mejor control. Con este objetivo establecido, comienza el desarrollo de ``Mezfer Insider'', un portal para administradores y clientes que unifica las funcionalidades de las aplicaciones anteriores en un solo lugar.
    \chapter{Objetivos}
    \section{Objetivo General}
    \begin{itemize}
        \item Desarrollar e implementar la aplicación ``Mezfer Insider'' para centralizar y controlar las actividades de administradores y clientes.
    \end{itemize}
    \section{Objetivos específicos}
\begin{itemize}
    \item Analizar la base de datos anterior para identificar áreas de mejora.
    \item Diseñar la nueva base de datos para Mezfer Insider.
    \item Integrar Mezfer Insider a Control Manager Siso (CMS) para que los administradores puedan iniciar sesión.
    \item Crear el inicio de sesión para los clientes de Mezfer.
    \item Crear el dashboard para los clientes.
    \item Integrar la funcionalidad de canje de puntos para los clientes.
\end{itemize}
    \chapter{Problemas a resolver}
    \begin{itemize}
        \item Fragmentación de plataformas: Actualmente, el programa Mezfer Rewards es gestionado mediante dos aplicaciones separadas: una aplicación web para los administradores y una aplicación móvil para los clientes. Esta separación implica un manejo dividido e ineficiente de las funcionalidades.
        \item Descentralización: El hecho de que las actividades estén dispersas en diferentes plataformas dificulta la capacidad de mantener una visión centralizada y controlada de las operaciones.
        \item Falta de integración: Las actividades de los clientes y los administradoes están divididas en distintos sistemas, lo que hace que la gestión de tareas sea poco eficiente.
        \item Falta de adaptación de funcionalidades: Se requiere unificar las funcionalidades de las aplicaciones web y móvil en una misma plataforma. Es importante que estas funcionalidades mantengan las características clave que ofrecen, como la gestión del programa y el acceso a las recompensas.
        \item Mejora de la experiencia del usuario: El usuario es una pieza fundamental en esta nueva plataforma, por lo que es importante asegurar que tanto administradores como clientes puedan realizar sus actividades eficientemente y sin problemas para ofrecer una experiencia agradable.    
    \end{itemize}
\end{document}